%! Convolutional Neural Nets — Unit 8, Lesson 8.2 — Student Packet Cover Sheet
\documentclass[10pt]{article}
\usepackage{linalg-article}
\usepackage{linalg-boxes}
\usepackage{ltablex}
\keepXColumns

\begin{document}

% Full-bleed burgundy banner
\begin{tikzpicture}[remember picture, overlay]
  \fill[burgundy] (current page.north west)
    rectangle ([yshift=-0.9in]current page.north east);
\end{tikzpicture}
\vspace{-0.8in}

\begin{center}
  {\color{white}\textbf{\LARGE Linear Algebra}}\\[4pt]
  {\color{blushmid}\large\textbf{Unit 8 \quad Learning from Data}}\\[6pt]
  {\color{white}\textbf{\large Lesson 8.2 \quad Convolutional Neural Nets}}
\end{center}

\vspace{0.22in}
\namedateperiod
\vspace{0.18in}

\begin{learningtargetbox}
\small
\begin{enumerate}[leftmargin=1.5em, itemsep=5pt]
  \item \ldots{} slide a small \textbf{filter} across a signal --- a dot product at each spot --- and read where it \textbf{fires}.
  \item \ldots{} write a convolution as a matrix whose weights \textbf{repeat down the diagonals}, and count how many weights \textbf{weight sharing} saves.
  \item \ldots{} explain how a convolutional layer $\mathrm{ReLU}(A\mathbf{v}+\mathbf{b})$ reuses one filter everywhere --- so it detects the same pattern \emph{wherever} it appears, with almost no weights.
\end{enumerate}
\end{learningtargetbox}

\vspace{0.14in}

\begin{tocbox}
\small
\renewcommand{\arraystretch}{1.5}
\begin{tabularx}{\linewidth}{@{} c l X r @{}}
\rowcolor{blush}
\textbf{\#} & \textbf{Component} & \textbf{Description} & \textbf{Score} \\
\midrule
1 & Warm-Up        & Spiral review: window dot products, a constant-diagonal matrix $\times$ vector, and ReLU & \blank{1.2cm} \\
2 & Guided Notes   & Sliding a filter; the feature map; the convolution matrix; weight sharing; a conv layer & \blank{1.2cm} \\
3 & Group Activity & Edge vs.\ smoothing filters, building the matrix, counting weights, detecting a pattern anywhere & \blank{1.2cm} \\
4 & Exit Ticket    & Quick check: slide a filter, count shared weights, and why one filter finds a pattern anywhere & \blank{1.2cm} \\
5 & Homework       & Sliding filters, the convolution matrix, weight counts, and the bias as a threshold & \blank{1.2cm} \\
\midrule
  & \multicolumn{2}{r}{\textbf{Total}} & \blank{1.2cm} \\
\end{tabularx}
\end{tocbox}

\vspace{0.10in}

\begin{remindbox}
\small A \textbf{full} layer gives every connection its own weight --- far too many for an image, and it
relearns the same pattern at every spot. A \textbf{convolution} fixes both: take a small \textbf{filter} (a
few weights) and \textbf{slide} it across the input, taking a dot product at each position. As a matrix, the
filter's weights \textbf{repeat down the diagonals} --- the same weights \emph{reused everywhere}
(\textbf{weight sharing}). So one small filter detects its pattern \emph{wherever} it appears, and a
convolutional layer $\mathrm{ReLU}(A\mathbf{v}+\mathbf{b})$ needs a handful of weights instead of a trillion.
\end{remindbox}

\end{document}
